\cotefile{36}
\scanpage{10}
\handwritten{\unsure{I|1}) \uline{Le théorème de Griffiths par voie algébrique.}}

1.\quad Soit
\[
f:Y\handwritten{\longrightarrow}S
\]
un morphisme propre et lisse de schémas, avec $S$ \inserted{\handwritten{régulier,}} connexe à caractéristiques résiduelles toutes $\ne\unsure{\chi|\ell}$ (nombre premier fixé). Soit $s\handwritten{\in}S$, $\bar s$ un point géométrique au dessus de $s$. Soi\deleted{\unsure{t}}\handwritten{t} $F$ un faisceau $\chi$-adique constant tordu sur $S$ (en pratique, un faisceau de la forme $\underline{Z}_{\chi}(i)$ ), et $x$ une classe de cohomologie
\[
x\handwritten{\in}H^{\deleted{N}\handwritten{m}}(Y,\deleted{\unsure{f(F))|f^*(F))}}\handwritten{F_Y)}\ .
\]
Disons que $x$ est \uline{primitive} si son image dans $H^0(S,R^m f_{\handwritten{*}}\handwritten{(\unsure{F|f}))}\deleted{\unsure{(F))|(F)}}$ est nulle, ou ce qui revient au même par le th. de spécialisation de la cohomologie, si son image dans $H^{\deleted{N}\handwritten{m}}(Y_{\bar s},F_{\bar s})$ est nulle. Alors la suite spectrale de Leray pour $f$ nous donne un élément
\[
u(x)\handwritten{\in}E_2^{1,\deleted{\unsure{N|N-1|m-1}}\handwritten{m-1}}=H^1(S,R^{\deleted{\unsure{m|N}}\handwritten{m-1}}_{\deleted{\unsure{2|f}}}f_{\deleted{\unsure{f}}\handwritten{*}}(F_{\handwritten{Y}}))\ .
\]
Donnons une propriété fonctorielle générale de cette classe $u(x)$. Pour ceci, soit $Y'$ lisse et propre sur $S$, et
\[
\unsure{Z|\mathbb{Z}|z}\handwritten{\in}H^N(Y\mathbin{\mathrm{x}}_S Y',G),
\]
de sorte que par le formalisme des correspondances cohomologiques, $z$ définit une application
\[
x'\handwritten{\longmapsto}Z(x'): \quad H^{m'}(Y',F')\handwritten{\longrightarrow}H^m(Y,F),
\]
si on suppose $m=m'+N-\deleted{\unsure{\dim(Y'/S)|\dim(X'/S)}}\handwritten{\unsure{2d'}}$, et qu'on a un accouplement $F'\mathbin{\ooalign{\hfil$\mathrm{x}$\hfil\cr\hfil$\handwritten{\bigcirc}$\hfil\cr}}G\handwritten{\longrightarrow}F(d')$ (NB $d'$ est la dimension relative de $Y'$ sur $S$, supposée constante, et $(d')$ désigne le twist de Tate\deleted{\unsure{xx|de}}; $F'$ et $G$ désignent des faisceaux $\chi$-adiques constants tordus \inserted{\handwritten{sur $S$}} tout comme $F$). Ceci dit, si $x'$ est primitive, alors il en est évidemment de même de $x=Z(x')$, et on aura de plus
\[
u(Z(x'))=Z(u(x'))\ ,
\]
où le deuxième membre s'explicite en notant que $Z$ définit également, par
% archive sheet number: 9
% The prime-number glyph is mechanically chi-shaped. The first occurrence retains
% chi/ell ambiguity; later chi follows the first candidate, not a mathematical
% assertion that the author intended the Greek letter.
% The pairing sign is a typed x with a thin handwritten oval around it; the
% standard-LaTeX overlay above records the components without adding a macro.
% Faint disconnected traces at the extreme upper-left edge are not securely
% identifiable as authorial writing; see report and inventory item 010-54.
